\section{Retaining the selected data through the final construction}
\label{sec:public-run-binding}

We finish the selection argument by specifying which common data are
retained. The issue is the order of the choices: the local estimates,
the numerical purchases and the limiting observables must all refer to
one family. The analytic hypotheses used here are those of
Sections~\ref{sec:common-funded-history},
\ref{sec:common-observable-purchases} and
\ref{sec:common-scale-selection}. This argument identifies their common
application; it supplies no replacement for an outstanding local estimate.

Fix a simple gauge group $G$ and a positive kernel penalty $b$.
Use the profile construction with support exponent exactly $100$ at
every stage, derivative order at least $j$ and spatial power at least
$j$ at stage $j$. Denote the selected kernel by $K$, its stage profiles
by $\pi_j$, and its primary and secondary row-price families by $P$ and
$P'$. The flat-step, curvature-quadratic symmetry and dual-support estimates select their
floors from these data before the block lengths are chosen. Take their
maximum together with the original, comparison, treasury and period
floors, and call it $n_j$. Raising the floor to $n_j$ changes neither
$K$ nor $\pi_j$ nor $P$. Select block lengths $L_j\ge n_j$ and retain
all these data when applying the scheduling construction.

Let $J(m)$ be the resulting stage clock, $\widehat\beta_m$ its hatted
coupling, and $H$ its quadratic horizon. Its actual block at step $m$
has length $L_{J(m+1)}$. For clarity, the two kinds of numerical payment
are different. If $U_{E,j}>0$ is the already selected width at fixed
index $E$, the countable purchase is
\[
 \forall E\ \exists m_E\ \forall m\ge m_E\ \forall j\le J(m+1):
 \quad U_{E,j}^{-1}\le 1+\log\log\widehat\beta_m.
\]
The same construction keeps $1\le1+\log\log\widehat\beta_m$ for
all $m$. After the lengths are fixed, let $A_j(L)$ and $B_j(L)$
be the two actual post-length costs and set
$F_j(L)=\max\{A_j(L),B_j(L)\}$. The all-stage purchase is
\begin{equation}\label{eq:public-native-allstage}
       F_{J(m+1)}(L)\le\widehat\beta_m\qquad(m\ge0).
\end{equation}
The countable version is retained as well. In particular, an eventual
statement for fixed $E$ is not being substituted for the all-stage
payment at $m=0$.

Here is the verification that the local suppliers apply to $H$ itself.
The selected levels have kernel $K$ and stage profiles $\pi_j$; the
horizon satisfies the row prices $P$. At step $m$, each component floor
is at most $n_{J(m+1)}\le L_{J(m+1)}$. The actual patch and the patch
used in computing $F_{J(m+1)}(L)$ agree, because a patch is determined
by its stage and block length. Thus \eqref{eq:public-native-allstage}
pays both actual costs. In the cost that has the form $e^{a_j}$, its
payment gives $a_j\le\log\widehat\beta_m$ by monotonicity of the
logarithm. The other payment supplies the flat-step delay. The flat,
dual-support and curvature-quadratic symmetry suppliers therefore apply with the original
levels and $H$, at every allowed cutoff and level. No new horizon is
chosen in this application.

Now retain the finite-volume history, its campaign and the transition
witnesses obtained from this one construction. Denote them by
$R$, $T$ and $W$, respectively. Here $R$ names the full selected
history and scale data, $T$ names its campaign of parameter paths,
and $W$ keeps the transition and finite-volume witnesses for that
campaign. The family $C_R$ is the family of
Section~\ref{sec:original-history-portfolio-attachment}. In addition
to the fixed penalty $b$, keep the treasury and common-period estimates
and the order and power inequalities when making this simultaneous
choice. They are conclusions of the selection before the history is
packaged. One may retain all of them with the same existential witness;
forgetting them in a shorter statement does not prove that they can be
recovered from its remaining hypotheses.

There are also two distinct support-exponent facts. The stage
parameters used by the operator-product construction have exponent
exactly $100$. The larger budget used to bound local functions has
support exponent at least $100$. The first fact follows from the fixed
profile under reindexing by $J$; the second is a budget bound. Neither
requires changing the selected history or replacing its profile by an
increasing support-exponent profile.

The funded continuation is now evaluated at the history, reference
schedule, local sections, transition data and finite-volume witnesses
stored in $R$. Proposition~\ref{prop:ex-original-history-attachment}
identifies its output family with $C_R$. Its trace and coverage
conclusions consequently concern that same family. For the
operator-product estimates, keep $W$ as well as $R,T$: the finite
transition witnesses enter the service of
Section~\ref{sec:ope-history-service}. The fixed exponent, common
clock, halo and logarithmic estimates are those just selected.
Thus the common-history, observable and operator-product applications
are simultaneous applications of their canonical arguments.

Finally this binding persists through the ordered extractions. Given
an initial extraction order $f$, first refine it to obtain the coupling
limits, then refine again to obtain the field limits. If the two
refinements are $g,w$, their combined order is
$f\circ g\circ w$, with the corresponding ordered period refinements.
Every coupling limit survives the second refinement, by the defining
limit property of a refinement. The phase, trace and operator-product
constructions are applied to the resulting single limiting table and
its single reconstructed system. Full-mesh convergence additionally
uses the full-mesh estimates on $C_R$; it is not inferred from the
existence of a mesh subsequence. These operations preserve the selected
history and its purchases while making the further choices required
by the limiting construction.
